\begin{proof}[Proof of \cref{obj:thm:degree-frontier}]
\begin{enumerate}
\item Apply \cref{obj:thm:bounded-outcome-degree-frontier} with the fixed \(\beta\) and \(p\).  It gives constants \(a,b\) with \(0<a\le b\) and, for every admissible \(n,d,B\), the block-energy lower bound, the projected-SNIPE upper bound, the unprojected-SNIPE upper bound, the local-energy comparison, the exposed-order comparison
\[
a\,d\binom{d}{k_\star(d,\beta,p)}
\le dA_d
\le b\,d\binom{d}{k_\star(d,\beta,p)},
\]
and the block-mixture conclusions.  Set
\[
c_{\beta,p}:=\min\{1,a^2\},
\qquad
C_{\beta,p}:=\max\{1,b^2\}.
\]
Then \(0<c_{\beta,p}\le C_{\beta,p}\). % lean: degree_frontier

\item Fix \(n,d\in\mathbb N\) and \(B\in\mathbb R\) with \(n\ge1\), \(d\le n\), and \(B\ge0\), and write
\[
x:=\frac{d\binom{d}{k_\star(d,\beta,p)}}{n},
\qquad
y:=\frac{dA_d}{n}.
\]
The common-\(p\) product Bernoulli design satisfies \cref{obj:ass:bernoulli-design}.  Since \(n\ge1\), the exposed-order comparison from the preceding item gives
\[
a x\le y\le b x,
\]
and hence
\[
\min\{1,ax\}\le \min\{1,y\}\le \min\{1,bx\}.
\]
Also \(\min\{1,a^2\}\le a\): if \(a\le1\), then \(\min\{1,a^2\}\le a^2\le a\), while if \(a>1\), then \(\min\{1,a^2\}\le1\le a\).  Splitting on \(x\le1\) and \(ax\le1\), and using this comparison in the saturated branches, gives
\[
\min\{1,a^2\}\min\{1,x\}
\le
a\min\{1,ax\}.
\]
Similarly \(b\le\max\{1,b^2\}\): if \(b\le1\), then \(b\le1\le\max\{1,b^2\}\), while if \(b>1\), then \(b\le b^2\le\max\{1,b^2\}\).  Splitting on \(x\le1\) and \(bx\le1\) gives
\[
b\min\{1,bx\}
\le
\max\{1,b^2\}\min\{1,x\}.
\]
In the branch \(x>1\) and \(bx\le1\), the needed bound is obtained by first deriving \(b<1\): otherwise \(b\ge1\) would imply \(x\le bx\), contradicting \(x>1\) and \(bx\le1\). % lean: degree_frontier

\item The lower bound follows by chaining the previous inequalities with the block-energy lower bound supplied by \cref{obj:thm:bounded-outcome-degree-frontier} and the identification of the coefficient-mass minimax risk in \cref{obj:def:minimax-risk}:
\[
\begin{aligned}
c_{\beta,p}B^2\min\{1,x\}
&=B^2\min\{1,a^2\}\min\{1,x\} \\
&\le B^2a\min\{1,ax\} \\
&\le B^2a\min\{1,y\} \\
&\le R_n^\star(d,\beta,p,B).
\end{aligned}
\] % lean: degree_frontier

\item The projected-SNIPE upper bound at the block-energy scale from \cref{obj:thm:bounded-outcome-degree-frontier}, together with the upper rescaling above, gives
\[
\begin{aligned}
\sup_{(G,c)\in\mathcal M_{n,d,\beta}(B)}
\mathbb E\!\left[
  \bigl(\widehat\tau_n^{\mathrm{up}}-\tau_n\bigr)^2
\right]
&\le bB^2\min\{1,y\} \\
&\le bB^2\min\{1,bx\} \\
&\le C_{\beta,p}B^2\min\{1,x\}.
\end{aligned}
\] % lean: degree_frontier

\item The estimator \(\widehat\tau_n^{\mathrm{up}}\) is an admissible competitor for the minimax problem.  The finite-population estimator is measurable, and \cref{obj:ass:bounded-coefficient-mass} gives, for every model in \(\mathcal M_{n,d,\beta}(B)\),
\[
|Y_i(z)|\le
\sum_{S\subseteq N_i}|c_{i,S}|\le B.
\]
Therefore
\[
|\tau_n|
\le
\frac1n\sum_{i\in V_n}\bigl(|Y_i(\mathbf 1)|+|Y_i(\mathbf 0)|\bigr)
\le 2B.
\]
The projection defining \(\widehat\tau_n^{\mathrm{up}}\) places the estimator in \([-B,B]\), so its pointwise squared loss is bounded by \(9B^2\), and hence its risk is uniformly bounded over the class.  Since \(n\ge1\), the empty graph with the zero coefficient schedule supplies a member of \(\mathcal M_{n,d,\beta}(B)\).  The definition of \(R_n^\star(d,\beta,p,B)\) as the infimum over admissible estimators gives
\[
R_n^\star(d,\beta,p,B)
\le
\sup_{(G,c)\in\mathcal M_{n,d,\beta}(B)}
\mathbb E\!\left[
  \bigl(\widehat\tau_n^{\mathrm{up}}-\tau_n\bigr)^2
\right].
\] % lean: snipeClipped_admissible, minimaxRisk_le_worstRisk_of_admissible

\item The unprojected SNIPE bound is obtained from the block-energy bound in \cref{obj:thm:bounded-outcome-degree-frontier}:
\[
\begin{aligned}
\sup_{(G,c)\in\mathcal M_{n,d,\beta}(B)}
\mathbb E\!\left[
  \bigl(\widehat\tau_n^{\mathrm{SNIPE}}-\tau_n\bigr)^2
\right]
&\le B^2\,\frac{dA_d}{n} \\
&\le B^2\,b\,x \\
&\le C_{\beta,p}B^2\,x.
\end{aligned}
\] % lean: degree_frontier

\item The score-energy assertion is exactly the local-energy comparison supplied by \cref{obj:thm:bounded-outcome-degree-frontier}: for every \((G,c)\in\mathcal M_{n,d,\beta}(B)\) and every \(i\in V_n\),
\[
A_i\le A_d.
\] % lean: degree_frontier

\item For every \((G,c)\in\mathcal M_{n,d,\beta}(B)\), every \(i\in V_n\), and every order \(r\) with \(1\le r\), the overlap-count assertion is \cref{obj:lem:overlap-count} applied to the graph \(G\), the degree bound \(d\), the order \(r\), and the unit \(i\):
\[
\sum_{\ell\in V_n}\binom{|N_i\cap N_\ell|}{r}
=
\sum_{\substack{S\subseteq N_i\\ |S|=r}}
\#\{\ell\in V_n:S\subseteq N_\ell\}
\le
d\binom{|N_i|}{r}
\le
d\binom{d}{r}.
\] % lean: overlap_count_le

\item Assume \(B>0\) and \(d\ge1\).  The block-mixture component of \cref{obj:thm:bounded-outcome-degree-frontier}, with \(m=\lfloor n/d\rfloor\), \(\rho\) the active population fraction, and \(\delta\) the least-favourable perturbation amplitude, gives the two complete-block models \(M_+\) and \(M_-\) for every \(U=(U_1,\ldots,U_m)\) with \(|U_b|\le B/2\).  Their graph is the complete-block graph, their coefficient schedules are the corresponding \(+1\) and \(-1\) block schedules, and
\[
\tau_n(M_+)=\rho\delta,
\qquad
\tau_n(M_-)=-\rho\delta.
\]
The same component gives
\[
H^2(\Pi_+,\Pi_-)
\le
\frac{4\pi^2m\delta^2}{B^2A_d},
\qquad
\frac12\le\rho\le1,
\qquad
\frac{A_d}{m}=\frac{dA_d/n}{\rho}.
\]
This proves the block-construction assertion and completes the proof. % lean: degree_frontier
\end{enumerate}
\end{proof}
